% ============================================================
% Sdílená preambule pro český akademický / technický TeX dokument
% Engine: XeLaTeX (vyžadováno kvůli fontspec/polyglossia)
% ============================================================

% --- Jazyk a fonty ---
\usepackage{fontspec}
\usepackage{polyglossia}
\setmainlanguage{czech}
% BasicTeX nemá fonty registrované v macOS — načítáme přímo soubory přes kpathsea
\setmainfont{lmroman10-regular.otf}[
  BoldFont = lmroman10-bold.otf,
  ItalicFont = lmroman10-italic.otf,
  BoldItalicFont = lmroman10-bolditalic.otf,
  Ligatures = TeX ]
\setsansfont{lmsans10-regular.otf}[
  BoldFont = lmsans10-bold.otf,
  ItalicFont = lmsans10-oblique.otf,
  BoldItalicFont = lmsans10-boldoblique.otf,
  Ligatures = TeX ]
\setmonofont{lmmonolt10-regular.otf}[
  BoldFont = lmmonolt10-bold.otf,
  ItalicFont = lmmonolt10-oblique.otf ]

% --- Matematika ---
\usepackage{amsmath}
\usepackage{amssymb}
\usepackage{mathtools}

% --- Rozvržení stránky ---
\usepackage[a4paper, margin=2.5cm]{geometry}
\usepackage{parskip}
\usepackage{fancyhdr}
\pagestyle{fancy}
\fancyhf{}
\fancyhead[LE]{\small\leftmark}
\fancyhead[RO]{\small\rightmark}
\fancyfoot[C]{\thepage}
\renewcommand{\headrulewidth}{0.4pt}

% --- Barvy, grafika, boxy ---
\usepackage{xcolor}
\usepackage{tikz}
\usetikzlibrary{shapes.geometric, arrows.meta, positioning, calc, fit}
\usepackage[skins, breakable, theorems]{tcolorbox}
\usepackage{enumitem}
\usepackage[autostyle]{csquotes}
\usepackage{booktabs}
\usepackage{multirow}
\usepackage{tabularx}
\usepackage{longtable}
\usepackage{needspace}

% --- Barevná paleta ---
\definecolor{pojemModra}{RGB}{21, 101, 192}
\definecolor{prikladSeda}{RGB}{84, 110, 122}
\definecolor{otazkaOranzova}{RGB}{230, 126, 34}
\definecolor{reseniZelena}{RGB}{46, 125, 50}
\definecolor{pozorCervena}{RGB}{198, 40, 40}
\definecolor{kucharkaFialova}{RGB}{106, 27, 154}

% ============================================================
% Makra a prostředí (jednotný styl pro všechny kapitoly)
% ============================================================

% Anglický ekvivalent pojmu — používat při PRVNÍM výskytu každého pojmu
% Použití: úrok\en{interest}
\newcommand{\en}[1]{ \textit{(angl. #1)}}

% Seznam symbolů po vzorci
% Použití: \begin{kde} \item $FV$ — budoucí hodnota \en{future value} \end{kde}
\newenvironment{kde}%
  {\par\noindent\textit{kde:}\begin{itemize}[nosep, leftmargin=2.5em]}%
  {\end{itemize}\medskip}

% --- Pojem: definice s povinným anglickým ekvivalentem ---
% Použití: \begin{pojem}{Úrok}{interest} ... \end{pojem}
\newtcolorbox[auto counter, number within=chapter]{pojem}[2]{%
  breakable, enhanced,
  colback=pojemModra!4, colframe=pojemModra!4,
  borderline west={3pt}{0pt}{pojemModra},
  fonttitle=\bfseries, coltitle=pojemModra,
  title={Pojem~\thetcbcounter: #1 \textmd{\textit{(angl. #2)}}},
  attach title to upper={\par},
  boxrule=0pt, frame hidden, sharp corners,
  left=8pt, right=8pt, top=6pt, bottom=6pt
}

% --- Příklad: řešený příklad krok za krokem ---
% Použití: \begin{priklad}{Reálná úroková míra} ... \end{priklad}
\newtcolorbox[auto counter, number within=chapter]{priklad}[1]{%
  breakable, enhanced,
  colback=black!3, colframe=prikladSeda,
  fonttitle=\bfseries,
  title={Příklad~\thetcbcounter: #1},
  boxrule=0.8pt, arc=2pt,
  left=8pt, right=8pt, top=6pt, bottom=6pt
}

% --- Testová otázka: PŘESNÉ znění otázky ze zdroje ---
% Použití: \begin{otazka}{PF_tema_4.pdf} text otázky
%   \begin{enumerate}[label=\alph*)] \item ... \end{enumerate}
% \end{otazka}
\newtcolorbox[auto counter, number within=chapter]{otazka}[1]{%
  breakable, enhanced,
  colback=otazkaOranzova!5, colframe=otazkaOranzova,
  fonttitle=\bfseries,
  title={Testová otázka~\thetcbcounter\ \ \textmd{\small(zdroj: #1)}},
  boxrule=0.8pt, arc=2pt,
  left=8pt, right=8pt, top=6pt, bottom=6pt
}

% --- Řešení: správná odpověď + proč; proč jsou distraktory špatně ---
% Použití: \begin{reseni}{c} ... \end{reseni}   (argument = písmeno správné odpovědi)
\newtcolorbox{reseni}[1]{%
  breakable, enhanced,
  colback=reseniZelena!5, colframe=reseniZelena,
  fonttitle=\bfseries,
  title={Řešení: správná odpověď \MakeUppercase{#1})},
  boxrule=0.8pt, arc=2pt,
  left=8pt, right=8pt, top=6pt, bottom=6pt
}

% --- Odpověď: modelová odpověď na otevřenou (kontrolní) otázku ---
% Použití: \begin{odpoved} ... \end{odpoved}
\newtcolorbox{odpoved}{%
  breakable, enhanced,
  colback=reseniZelena!5, colframe=reseniZelena,
  fonttitle=\bfseries,
  title={Modelová odpověď},
  boxrule=0.8pt, arc=2pt,
  left=8pt, right=8pt, top=6pt, bottom=6pt
}

% --- Pozor: chytáky a časté chyby ---
\newtcolorbox{pozor}{%
  breakable, enhanced,
  colback=pozorCervena!5, colframe=pozorCervena,
  fonttitle=\bfseries,
  title={Pozor — častá chyba!},
  boxrule=0.8pt, arc=2pt,
  left=8pt, right=8pt, top=6pt, bottom=6pt
}

% --- Kuchařka: postup krok za krokem (čtení tabulek, integrační metody…) ---
% Použití: \begin{kucharka}{Jak číst tabulku Studentova rozdělení} \begin{enumerate} ... \end{enumerate} \end{kucharka}
\newtcolorbox[auto counter, number within=chapter]{kucharka}[1]{%
  breakable, enhanced,
  colback=kucharkaFialova!4, colframe=kucharkaFialova,
  fonttitle=\bfseries,
  title={Kuchařka~\thetcbcounter: #1},
  boxrule=0.8pt, arc=2pt,
  left=8pt, right=8pt, top=6pt, bottom=6pt
}

% --- Styly uzlů pro rozhodovací stromy (TikZ) ---
\tikzset{
  rozhodnuti/.style={diamond, aspect=2.2, draw=kucharkaFialova, fill=kucharkaFialova!8,
                     text width=9.2em, align=center, font=\small, inner sep=1pt},
  vysledek/.style={rectangle, rounded corners=3pt, draw=reseniZelena, fill=reseniZelena!8,
                   text width=10em, align=center, font=\small, inner sep=4pt},
  sipka/.style={-{Stealth}, thick},
  popisek/.style={font=\footnotesize\itshape, midway}
}

% --- Věta: matematická věta / tvrzení (rigorózní kurzy) ---
% Použití: \begin{veta}{Lagrangeova věta} ... \end{veta}
\definecolor{vetaTeal}{RGB}{0, 121, 107}
\newtcolorbox[auto counter, number within=chapter]{veta}[1]{%
  breakable, enhanced,
  colback=vetaTeal!4, colframe=vetaTeal,
  fonttitle=\bfseries, coltitle=vetaTeal,
  title={Věta~\thetcbcounter: #1},
  attach title to upper={\par},
  boxrule=0pt, frame hidden, sharp corners,
  borderline west={3pt}{0pt}{vetaTeal},
  left=8pt, right=8pt, top=6pt, bottom=6pt
}

% --- Důkaz: jednoduché prostředí s QED značkou ---
% Použití: \begin{dukaz} ... \end{dukaz}
\newenvironment{dukaz}%
  {\par\medskip\noindent\textit{Důkaz.}\ }%
  {\hfill$\square$\par\medskip}

% --- Hyperref (vždy poslední) ---
\usepackage[unicode, colorlinks=true, linkcolor=pojemModra, urlcolor=pojemModra]{hyperref}
